\documentclass[11pt]{article}

\usepackage[a4paper,margin=2.5cm]{geometry}
\usepackage{amsmath}
\usepackage{booktabs}
\usepackage{listings}
\usepackage{xcolor}
\usepackage{hyperref}

\lstset{
  basicstyle=\ttfamily\small,
  breaklines=true,
  columns=fullflexible,
  frame=single,
  backgroundcolor=\color{gray!6}
}

\title{BernoulliSolver CLI Validation Testcase 03}
\author{}
\date{}

\begin{document}
\maketitle

\section{Purpose}

This report documents the execution of CLI validation testcase 03 from
\texttt{Bending\_Torsion\_Beam\_Test\_Case/BernoulliSolver.py}. The case is a
simply supported beam under a uniform distributed load represented by
equivalent nodal forces and end moments.

The terminal output of the executed run was saved to:

\begin{lstlisting}
Bending_Torsion_Beam_Test_Case/BernoulliSolver_python_TestCases/03_simply_supported_udl/cli_validation_03_terminal_output.txt
\end{lstlisting}

\section{Input Model}

The input model is:

\begin{lstlisting}
Bending_Torsion_Beam_Test_Case/TestCase_Input/03_simply_supported_udl/simply_supported_udl.mdpa
\end{lstlisting}

The model contains 11 nodes and 10 two-node beam elements. The beam lies on the
global $x$ axis:

\begin{lstlisting}
Node 1:  x = 0
Node 11: x = 10
Element length: 1
\end{lstlisting}

\section{Boundary Conditions And Loads}

The supports are prescribed through command-line constraints:

\begin{lstlisting}
Node 1:
  ux = 0
  uy = 0
  uz = 0
  rx = 0

Node 11:
  uy = 0
  uz = 0
  rx = 0
\end{lstlisting}

The uniform distributed load is:

\begin{equation}
q_y = -1000 .
\end{equation}

It is represented by equivalent nodal loads:

\begin{lstlisting}
Node 1:  Fy = -500,  Mz = -83.3333333333
Nodes 2-10: Fy = -1000
Node 11: Fy = -500, Mz = +83.3333333333
\end{lstlisting}

\section{Command Used}

The following command was run from the Kratos root directory:

\begin{lstlisting}[language=bash]
python3 Bending_Torsion_Beam_Test_Case/BernoulliSolver.py \
  Bending_Torsion_Beam_Test_Case/TestCase_Input/03_simply_supported_udl/simply_supported_udl.mdpa \
  --materials Bending_Torsion_Beam_Test_Case/Bending_Torsion_Beam_Test_Case/beam_geometry/StructuralMaterials.json \
  --prescribe 1 ux 0 --prescribe 1 uy 0 --prescribe 1 uz 0 --prescribe 1 rx 0 \
  --prescribe 11 uy 0 --prescribe 11 uz 0 --prescribe 11 rx 0 \
  --load 1 0 -500 0 0 0 -83.3333333333 \
  --load 2 0 -1000 0 0 0 0 \
  --load 3 0 -1000 0 0 0 0 \
  --load 4 0 -1000 0 0 0 0 \
  --load 5 0 -1000 0 0 0 0 \
  --load 6 0 -1000 0 0 0 0 \
  --load 7 0 -1000 0 0 0 0 \
  --load 8 0 -1000 0 0 0 0 \
  --load 9 0 -1000 0 0 0 0 \
  --load 10 0 -1000 0 0 0 0 \
  --load 11 0 -500 0 0 0 83.3333333333 \
  --output Bending_Torsion_Beam_Test_Case/BernoulliSolver_python_TestCases/03_simply_supported_udl
\end{lstlisting}

\section{Runtime Summary}

The solver printed:

\begin{lstlisting}
Solved reduced system with 59 free DOFs.
Read 11 nodes and 10 elements.
Wrote VTK output to: Bending_Torsion_Beam_Test_Case/BernoulliSolver_python_TestCases/03_simply_supported_udl
\end{lstlisting}

\section{Nodal Results}

The following table lists the nodal displacement in the load direction and the
corresponding bending rotation around $z$.

\begin{center}
\begin{tabular}{rrrr}
\toprule
Node & $x$ & $u_y$ & $\theta_z$ \\
\midrule
1  & 0  &  0.000000000000000e+00 & -2.013855324633510e-07 \\
2  & 1  & -1.975592073465477e-07 & -1.901079426454041e-07 \\
3  & 2  & -3.737715482519807e-07 & -1.594973417109750e-07 \\
4  & 3  & -5.117206379893772e-07 & -1.143869824391844e-07 \\
5  & 4  & -5.993233446109359e-07 & -5.961011760915261e-08 \\
6  & 5  & -6.293297889479755e-07 &  5.607720339794012e-23 \\
7  & 6  & -5.993233446109358e-07 &  5.961011760915284e-08 \\
8  & 7  & -5.117206379893771e-07 &  1.143869824391844e-07 \\
9  & 8  & -3.737715482519807e-07 &  1.594973417109747e-07 \\
10 & 9  & -1.975592073465479e-07 &  1.901079426454040e-07 \\
11 & 10 &  0.000000000000000e+00 &  2.013855324633515e-07 \\
\bottomrule
\end{tabular}
\end{center}

\section{Analytical Reference}

For a simply supported Euler-Bernoulli beam under a uniform distributed load,
the maximum displacement at midspan and the support rotation are:

\begin{equation}
u_y(L/2) = \frac{5q_yL^4}{384EI_z},
\end{equation}

\begin{equation}
\theta_z(0) = \frac{q_yL^3}{24EI_z}.
\end{equation}

Using:

\begin{lstlisting}
q_y = -1000
L   = 10
E   = 2.0690e11
I_z = 1
\end{lstlisting}

the analytical values are:

\begin{lstlisting}
u_y(L/2)      = -6.293297889479620e-07
theta_z(0)   = -2.013855324633478e-07
\end{lstlisting}

\section{Error Summary}

\begin{center}
\begin{tabular}{lrrrr}
\toprule
Quantity & Analytical & Python & Absolute error & Relative error \\
\midrule
$u_y(L/2)$
& -6.293297889479620e-07
& -6.293297889479755e-07
& 1.344664803766202e-20
& 2.136661615866065e-12\% \\
$\theta_z(0)$
& -2.013855324633478e-07
& -2.013855324633510e-07
& 3.202843331805323e-21
& 1.590403884841251e-12\% \\
\bottomrule
\end{tabular}
\end{center}

\section{Conclusion}

CLI validation testcase 03 ran successfully. The Python Bernoulli solver
matches the analytical midspan displacement and support rotation to
near-machine precision for this discretized equivalent-nodal-load setup. The
temporary terminal output block marked with
\texttt{TEMP DEBUG OUTPUT - REMOVE AFTER VERIFICATION} also printed all nodal
displacements and rotations directly to the terminal log.

\end{document}
